% Owner: Kais
\section{Message}\label{sec:message}

A beacon message must say \emph{what} was found, \emph{where} to go and
\emph{when} it was written, in a frame small enough for one ESP-NOW packet
(250 bytes). Fields are big-endian, and every frame ends with a checksum, the
XOR of all preceding bytes.

\paragraph{Identifier.} Every node that writes frames has a unique 1-byte
ID: \texttt{0x00} for the Command Post, \texttt{0x01}--\texttt{0x7F} for
Writers, \texttt{0x80}--\texttt{0xFE} for Executors. A frame is identified by
(frame type, \texttt{origin\_id}, \texttt{msg\_id}), where \texttt{msg\_id} is
a counter unique per origin. This tells which Writer composed each message,
and supports any number of robots. A beacon has no ID of its own: it is named
by the identifier of the frame it stores.

\paragraph{Event frame.} A Writer composes an EVENT\_BEACON frame and stores
it in the beacon it drops (Table~\ref{tab:event-frame}). It is $14 + 2N$
bytes long, where $N$ is the number of movement primitives (at most 118).

\begin{table}[h]
\centering
\small
\begin{tabularx}{\textwidth}{@{}l l l X@{}}
\toprule
Byte & Field & Size & Content \\
\midrule
0 & frame\_type & 1 & \texttt{0x01} \\
1 & origin\_id & 1 & Writer that composed it \\
2--3 & msg\_id & 2 & counter, unique per Writer \\
4 & ttl & 1 & hops left; decremented by each relay \\
5 & seq\_num & 1 & drop order in this Writer's run \\
6--9 & timestamp & 4 & seconds since this Writer's GO\_TRIGGER \\
10 & event\_type & 1 & \texttt{0x00} none (connectivity beacon), \texttt{0x01}
  gas, \texttt{0x02} fire, \texttt{0x03} victim alive, \texttt{0x04} victim
  dead, \texttt{0x05} robot down \\
11 & event\_value & 1 & severity, 0--255 \\
12 & primitive\_count & 1 & $N$ \\
13.. & primitives & $2N$ & path from the previous beacon: FORWARD (cm) or
  TURN (signed degrees) \\
last & checksum & 1 & XOR of all preceding bytes \\
\bottomrule
\end{tabularx}
\caption{EVENT\_BEACON frame.}
\label{tab:event-frame}
\end{table}

\paragraph{Where.} The primitives describe the path from the Writer's
previous beacon (from PEB-Inside for the first one). An Executor reaches an
event by following the beacons in drop order, and the Command Post rebuilds
every position by chaining these paths back to the entrance. Relative paths
keep frames short: their size depends on the distance between two beacons,
not on the distance from the entrance.

\paragraph{When, and message aging.} Messages age in three ways:
\begin{itemize}
  \item \textbf{Event age.} The timestamp counts seconds from the Writer's
    GO\_TRIGGER. The Command Post knows when it sent that order, so it turns
    every timestamp into real time and shows on the map how old each event is.
  \item \textbf{Hop limit.} Each relay decrements \texttt{ttl}, and a frame
    reaching 0 is no longer relayed, so no frame circulates forever.
  \item \textbf{Forgetting.} A beacon drops a frame it has already seen, but
    forgets it after 60\,s. Each frame is thus re-flooded about once a minute,
    and a node that arrives later, such as an Executor, still receives it.
\end{itemize}

\paragraph{Other frames.} Both start with the same first 5 bytes (type,
origin, \texttt{msg\_id}, \texttt{ttl}) and end with a checksum:
\begin{itemize}
  \item \textbf{GO\_TRIGGER} (8 bytes), from the Command Post: a
    \texttt{target\_id} (one robot, or all) and a \texttt{mission\_code}:
    explore or repair for a Writer, or the event type a robot must solve. Once the
    Executor is inside, PEB-Inside briefs it by replaying the stored event
    frames of that type, in drop order.
  \item \textbf{STATUS\_UPDATE} (8 bytes), from a beacon or a robot about
    itself: a status code (beacon low battery, or robot heartbeat) and its
    battery level (Section~\ref{sec:failure-cases}).
\end{itemize}
